O câncer de próstata é uma das neoplasias malignas mais comuns em homens, representando um importante desafio de saúde pública \cite{ref_siegel2024}. Estatísticas reportadas pela Agência Internacional para Pesquisa em Câncer indicam que, em 2022, essa neoplasia foi a quarta mais comum no mundo, correspondendo a 7,3\% de todos os casos de câncer \cite{JournalStatistic2024}. Nesse contexto, a classificação tumoral desempenha um papel central na definição do manejo clínico e no prognóstico. O sistema de Gleason, baseado na soma dos dois padrões histológicos predominantes, evoluiu para o sistema de classificação proposto pela Sociedade Internacional de Patologia Urológica (ISUP), organizado em cinco grupos variando de 1 a 5. Esse sistema fornece uma melhor estratificação prognóstica e é amplamente adotado na prática clínica \cite{ref-epstein2016}.

A avaliação histopatológica tradicional de imagens de lâmina inteira (Whole-Slide Images - WSIs) baseia-se na análise manual realizada por patologistas. Devido ao grande número de campos que devem ser examinados em uma única lâmina, o processo de avaliação torna-se uma tarefa repetitiva que exige longas jornadas de trabalho, atenção sustentada e conhecimento especializado. A fadiga decorrente da análise prolongada pode introduzir variabilidade inter e intraobservador nos resultados, influenciando diretamente o manejo clínico dos pacientes \cite{ref-bulten2020}. Consequentemente, nos últimos anos, tem havido um crescimento significativo no uso de técnicas de Inteligência Artificial (IA) para automatizar a análise de exames histopatológicos e auxiliar no diagnóstico do câncer de próstata, com o objetivo de mitigar essas limitações \cite{expert2022review}.

Nesse cenário, as Redes Neurais Convolucionais (CNNs) têm representado o estado da arte na análise de imagens médicas \cite{ref-campanella2019}, sendo capazes de extrair características relevantes e combinar padrões visuais complexos para identificar condições patológicas específicas \cite{araujo2021active}. Essa capacidade torna a abordagem promissora para a análise automatizada de padrões de Gleason e grupos de grau ISUP, reduzindo a variabilidade inerente à avaliação convencional. Com base nessa premissa, técnicas de \textit{deep learning} para classificação do câncer de próstata têm sido amplamente investigadas na literatura, com avanços relevantes em WSIs, estratégias de extração de patches e mecanismos de agregação multi-escala. Por exemplo, no contexto de aprendizado fracamente supervisionado, \cite{Xiang2023Prostate} propuseram o framework GCN-MIL, que utiliza uma Graph Neural Network para agregar características espaciais de patches em lâminas WSI. O modelo incorpora uma estratégia de treinamento robusto que filtra iterativamente amostras com alta incerteza para mitigar o impacto de rótulos ruidosos, alcançando um coeficiente kappa quadrático de 0.931 no conjunto interno do desafio PANDA. No entanto, a abordagem utiliza a função de perda tradicional \textit{Cross-Entropy}, que trata os graus ISUP como classes nominais independentes e ignora a hierarquia biológica da progressão do câncer.

Em uma validação clínica independente com 593 imagens do \textit{Seoul National University Hospital}, \cite{jung2022artificial} avaliaram o sistema DeepDx Prostate, alcançando um kappa quadrático de 0.922 na classificação de Gleason. Contudo, a generalização do modelo foi limitada pela homogeneidade do conjunto de dados institucional e por uma tendência de superestimar casos do Grupo de Grau 1 (GG1). Buscando a predição indireta de risco, \cite{liu2022deep} investigaram o uso de \textit{deep learning} para identificar pacientes com risco de câncer a partir de lâminas previamente classificadas como benignas. O método alcançou uma AUC de 0.739 para predição de risco, revelando padrões morfológicos sutis, mas foi limitado pela dependência de um único domínio institucional.

Para lidar com a variabilidade de escala das estruturas glandulares, \cite{pyramidalcnn} propuseram uma arquitetura \textit{Pyramidal CNN} composta por três redes rasas operando em diferentes tamanhos de patch. Embora o método tenha alcançado aproximadamente 0.77 de acurácia na classificação de padrões, foi limitado por rotulação semi-automática e desempenho moderado em classes mais complexas. Em paralelo, visando melhorar a interpretabilidade, \cite{kosoko2024xai_prostate} avaliaram arquiteturas combinadas com técnicas de \textit{Explainable AI} (\textit{Grad-CAM}) utilizando o conjunto SICAPv2. O modelo VGG19 alcançou o melhor desempenho (AUC-ROC de 0.85), embora o estudo não tenha incluído validação clínica externa. Seguindo uma estratégia de simplificação, \cite{afifi2024prostate} formularam o problema como uma tarefa de classificação binária (tecidos benignos vs. malignos). Utilizando \textit{fine-tuning} em modelos InceptionResNetV2 com a função de perda \textit{binary cross-entropy}, alcançaram 91,8\% de acurácia, enfatizando a detecção em detrimento da classificação multiclasse de Gleason.

Em arquiteturas mais avançadas, a incorporação de mecanismos de atenção também tem ganhado destaque. \cite{alici2024efficient} propuseram um modelo baseado na EfficientNet-B4 com mecanismo de atenção por canal (Eff4-Attn), alcançando 93,32\% de acurácia em uma tarefa de classificação com oito classes utilizando imagens do \textit{DiagSet}. No entanto, o estudo baseia-se exclusivamente em patches isolados e sofre com desbalanceamento de classes. Complementarmente, \cite{ikromjanov2021vit_prostate} investigaram o uso de \textit{Vision Transformers} (ViT) em um subconjunto do conjunto PANDA, alcançando acurácia superior a 85\% na classificação de patches, destacando a capacidade dos mecanismos de atenção em capturar dependências contextuais globais.

Apesar dos avanços recentes, diversas lacunas metodológicas ainda permanecem. Muitos estudos tratam o problema como uma tarefa de classificação nominal e utilizam funções de perda tradicionais, desconsiderando a natureza ordinal dos graus ISUP. Como consequência, erros de predição são penalizados de forma uniforme, sem considerar a magnitude da discrepância entre classes, o que não reflete adequadamente sua relevância clínica. Além disso, muitas abordagens não incorporam mecanismos robustos de filtragem de ruído ou modelagem de incerteza, aspectos importantes devido à presença de artefatos e inconsistências em conjuntos de dados públicos.

Para abordar essas limitações, este trabalho propõe um pipeline para a classificação automática dos grupos de grau ISUP em imagens histopatológicas. As principais contribuições são: (i) uma função de perda híbrida que combina penalização ordinal e \textit{focal loss}, considerando a natureza hierárquica do problema e o desbalanceamento de classes; (ii) uma estratégia de remoção de ruído baseada na incerteza estimada a partir da média das predições de modelos EfficientNet (B0, B1, B2 e B3); e (iii) uma avaliação comparativa das arquiteturas EfficientNet B0, B3 e B7, visando obter um modelo de inferência mais robusto e consistente.